\documentclass[uplatex,dvipdfmx,titlepage]{jsarticle}
\usepackage[dvipdfmx]{graphicx}
\usepackage{amsmath,amssymb}
\usepackage{bm}
\usepackage{enumitem} 

%レポート作成に関する注意
%レポートには実験の目的、内容、実験方法、実験結果、結果に関する吟味・考察、予想との乖離があるならその原因の考察を自分の言葉でまとめる

\title{固体力学の実験}
\author{京都大学工学部物理工学科 1025363386 伊藤温大}
\date{\today}

\begin{document}
\maketitle

% --------------------------------------------------
% 目次（不要な場合は以下の2行を消すか % でコメントアウトしてください）
\tableofcontents
\clearpage
% --------------------------------------------------

\section{はじめに}
本レポートでは超音波計測による弾性波伝搬速度と弾性係数の評価についてその結果、実験課題について報告する。


% ==================================================
% 実験2
% ==================================================
\section{超音波計測による弾性波伝搬速度と弾性係数の評価}
実験1\\
金属ブロックに縦波および、横波の超音波を送信、受診することによって測定波形から伝搬速度を求めた。
%この実験装置については形式を整える
この際、以下のような実験装置を使用した。
超音波送受信装置\\
トランスデューサー\\
音響結合媒質\\
試験片\\
具体的に試験片として以下の材料を用いた。
・機械構造用炭素鋼鋼材 S25C
・機械構造用炭素鋼鋼材 S45C
・アルミ二ウム合金     A5052
・アルミ二ウム合金     A2024(超ジュラルミン)
ディジタルオシロスコープ\\
パーソナルコンピューター\\
\subsection{実験方法}
1.試験片の寸法を測定し、記録した。
この時同様の材質について4回測定し、平均をとった。
記録した試験片の寸法は以下の図のようになった。
%有効数字について考察する
\begin{table}[]
\begin{tabular}{llll}
S25C  & S45C  & A5052 & A2025  \\
9.882 & 9.873 & 9.98  & 10.061
\end{tabular}
\end{table}
2.試験片を厚さ方向に複数回伝搬した波形について、それぞれの反射波の到着時刻を求めた。
この時S25Cの縦波について初めの波は超音波送受信装置から出た波を観測してしまっていると考えられる。
縦波について
受信したデータはS25Cから順に以下図のようになった。


この反射波の到着時刻を求める際には波の下限のピークの値を到着時刻として統一した。
この時縦波における到着時刻はS25Cから順に以下の図のようになった。





このとき
横波について
受診したデータは以下のようになった。






3.異なる反射回数(異なる伝搬距離)に対応する反射波の到達時刻を求めた。
4.受診した反射波形に含まれる周波数成分を、高速フーリエ変換により確認した。



実験2
1.送信トランスデューサと受信トランスデューサの間隔を変化させながら、波の到達時刻を調べた

2.トランスデューサ間の間隔と到達時刻との関係から伝搬速度を求めた。



考察課題
1.超音波測定に用いた4種類の材料に対して、ヤング率/密度として定義されている比剛性を求めると、、、
、、、、、、
2.超音波測定に用いた炭素鋼S25CとS45Cの違い、およびアルミニウム合金A5052とA2024の違いを図書で調べると、、、、、、、





\subsection{結果}

\subsection{考察}


% ==================================================
% まとめと参考文献
% ==================================================
\section{総括}
本実験を通して、光弾性法および超音波計測を用いた固体力学の評価手法について...

% 見出し番号をつけたくない場合は section に * をつけます
\section*{参考文献}
\begin{enumerate}[label={[\arabic*]}]
    \item 著者名, 『書籍名』, 出版社, 発行年.
    \item 物理工学科 実験指導書, 2026.
\end{enumerate}

\end{document}